\appendix{}

\chapter{Derivations}\label{app:derivations}

\section{Training-Time Sampling Laws}\label{app:training-time-laws}

The diffusion baseline draws a discrete denoising index uniformly,
$k\sim\operatorname{Uniform}\{0,\ldots,K-1\}$, where $K$ is the number of diffusion steps. This is not
a draw from the continuous transport time used by the flow-based models.

For instantaneous-velocity matching, let $U\sim\operatorname{Beta}(1.5,1)$ and set the transport
time to $\tau=1-U$, on the axis where $\tau=0$ is noise and $\tau=1$ is data. Equivalently,
\begin{equation}
  \tau\sim\operatorname{Beta}(1,1.5),\qquad
  p(\tau)=1.5(1-\tau)^{1/2},\quad 0<\tau<1.
\end{equation}
Thus the implemented draw weights the noise end of this transport axis. The expression
$\operatorname{Beta}(1.5,1)$ in the training configuration describes $U$ \emph{before} the reflection,
not the time passed to the network.

Analytic and consistency-interpolated average-velocity matching each draw two independent standard
normal variables $Z_1,Z_2$ and transform them by
\begin{equation}
  T_i=\frac{1}{1+\exp[-(P\sidx{std}Z_i-P\sidx{mean})]},\qquad
  P\sidx{mean}=-0.4,\quad P\sidx{std}=1.0,\quad i\in\{1,2\}.
\end{equation}
On the same noise-at-zero axis, $\operatorname{logit}(T_i)$ therefore has mean $+0.4$ and standard
deviation $1.0$. The ordered pair is $r=\min(T_1,T_2)$ and $t=\max(T_1,T_2)$. For each example,
independently with probability $0.5$, the implementation then sets $r=t$, giving a zero-length interval and an
instantaneous-velocity anchor; otherwise the interval has length $t-r>0$. Before ordering, the
logit-normal density of each $T_i$ is
\begin{equation}
  p_T(u)=\frac{1}{u(1-u)\sqrt{2\pi}}
  \exp\!\left[-\frac{(\operatorname{logit}(u)-0.4)^2}{2}\right],\qquad 0<u<1.
\end{equation}
These laws are the samplers behind the entries in \autoref{tab:train}; the average-velocity pair and
the losses to which it is supplied are defined in \autoref{sec:method:engine}.
\srcnote{Diffusion: \texttt{mix\_visual\_aligning/models/diffusion.py:299--301}. Reflected Beta:
\texttt{mix\_visual\_aligning/models/fm\_diffusion.py:300--307}. Logit-normal pair, ordering, and
zero-length fraction: \texttt{mix\_visual\_aligning/models/mf\_diffusion.py:353--368,423--428},
\texttt{mix\_visual\_aligning/models/af\_diffusion.py:486--505,691--696}; parameters in
\texttt{config/avoiding-d3il.py:895--897,969,1026--1028} and corresponding visual/UAV configs.}

\chapter{Extended Results}\label{app:results}

\section*{Intermediate diffusion budget on D3IL-avoiding}

\autoref{tab:avoiding-dpcc-protocol} locks the main comparison to one and two network evaluations,
with the diffusion model at its native twenty-step budget as the baseline of record. The available
intermediate diffusion checkpoint is reported here rather than inserted as an unmatched budget in that
comparison.

\begin{table}[htpb]
  \caption[Diffusion at ten denoising steps]{The diffusion model at $\nfe=10$ on D3IL-avoiding, at the
  protocol of \ac{DPCC}: five training seeds $\times$ two episodes per geometry, tightened constraints,
  four candidate plans and the temporal U-Net of $4.0$\,M parameters. Values are mean $\pm$ the standard
  deviation over the training seeds.}
  \label{tab:app:diffusion-k10}
  \centering
  \small
  \begin{tabularx}{0.82\linewidth}{@{}c rrr@{}}
    \toprule
    Rule & S\&C & Steps & ms/step \\
    \midrule
    $r$ & $0.900 \pm 0.091$ & $75.1 \pm 7.3$  & $327.2 \pm 35.4$ \\
    $c$ & $1.000 \pm 0.000$ & $70.3 \pm 11.5$ & $309.8 \pm 11.2$ \\
    $t$ & $1.000 \pm 0.000$ & $68.7 \pm 7.0$  & $321.7 \pm 14.4$ \\
    \bottomrule
  \end{tabularx}
\end{table}
\dataref{analysis\_results\_checkpoint/19-09-UAV-Pillars-Exclude/
batch\_avoiding\_combined\_20260919\_132703; DA\_in\_Paper/analysis/
avoiding\_rules\_by\_protocol.py}

\section{Twenty-Episode Evaluation on D3IL-avoiding}\label{app:avoiding-twenty-episode}

The extended evaluation uses twenty episodes per training seed and geometry, one hundred per geometry.
Coverage at $\nfe=20$ is listed in \autoref{tab:coverage-k20}.

What the larger sample buys is resolution. At ten episodes per geometry a success rate resolves to
$0.1$, so a rate of $1.00$ there is not evidence of a rate of $1.00$: the baseline's perfect rates fall
to $0.95$ and $0.92$ at one hundred episodes, while its tightened configurations move by at most $0.08$
(\autoref{tab:tiers}).

\begin{table}[htpb]
  \caption[Effect of the number of episodes]{The baseline at ten and at one hundred episodes per
  geometry: same configurations, same code.}
  \label{tab:tiers}
  \centering
  \begin{tabularx}{\linewidth}{@{}L cc@{}}
    \toprule
    Success with constraint satisfaction & 10 episodes & 100 episodes \\
    \midrule
    Cumulative projection cost, tightened, top-right-hard & 1.00 & 0.95 \\
    Temporal consistency, tightened, top-right-hard       & 1.00 & 0.92 \\
    Untightened configurations                            & \multicolumn{2}{c}{change by up to 0.35} \\
    Tightened configurations                              & \multicolumn{2}{c}{change by at most 0.08} \\
    \bottomrule
  \end{tabularx}
\end{table}

\dataref{DA\_20260819 --- baseline at 10 against 100 episodes per geometry}

At that resolution the selection rule separates the models more than the objective does
(\autoref{tab:state-headline}).

Analytic average-velocity matching is the model most sensitive to it. With temporal consistency it reaches
$0.993$ at one network evaluation, in $61.0$ control steps; with cumulative projection cost, the rule
\ac{DPCC} itself selects, the same model falls to $0.943$ and $72.0$ steps, and at two evaluations that rule
picks plans that stall, $98.0$ steps. Instantaneous-velocity matching is indifferent: every rule puts it
between $0.990$ and $1.000$ at one and two evaluations. The diffusion baseline does best with cumulative
projection cost, $0.983$ in $69.0$ steps.

Temporal consistency is therefore used for the flow-based operating points on this task. Under the
per-step projector the models differ, at their own best rule, by less than the gap in what one action
costs to compute.

The ordering is the one the protocol of \ac{DPCC} gives
(\autoref{tab:avoiding-dpcc-protocol}), and the two budgets do not change places: instantaneous-velocity
matching reads $0.997$, $0.997$ and $0.990$ at $\nfe=1$ against $0.960$, $0.970$ and $0.987$ at
$\nfe=20$ under the three rules, the largest gap between them being eleven episodes in three hundred.

\begin{table}[htpb]
  \caption[The generative models under every selection rule]{D3IL-avoiding at twenty episodes per seed:
  five training seeds $\times$ twenty episodes per geometry. Metrics, rules, marks and set-up as in
  \autoref{tab:avoiding-dpcc-protocol}; mean $\pm$ the sample standard deviation over the five seeds
  (\autoref{sec:setup:metrics:spread}). This protocol reaches the two budgets the flow-based models are
  operated at and the baseline's own budget of twenty; CI-MeanFM has it for one training seed and is
  reported as a single-seed matched-budget comparison in \autoref{tab:avoiding-budget20}.
  $^{\dagger}$ \textbf{The $\nfe=20$ rows do not rest on the full five seeds $\times$ three geometries
  the other rows do} --- the guard below counts what each of them covers. Their point estimate is the
  geometry-mean over the seeds present and their spread is over those seeds.}
  \label{tab:state-headline}
  \centering
  \footnotesize
  \begin{tabularx}{\linewidth}{@{}L r c rrr@{}}
    \toprule
    Model & $\nfe$ & Rule & S\&C & Steps & ms/step \\
    \midrule
    MeanFM     & 1 & $r$ & $0.930 \pm 0.032$ & $64.7 \pm 1.8$ & $18.4 \pm 0.1$ \\
                     &    & $c$ & $0.943 \pm 0.040$ & $72.0 \pm 4.1$ & $18.0 \pm 0.4$ \\
                     &    & \mainconfig{\selectedmark\ $t$} & \mainconfig{$0.993 \pm 0.015$} & \mainconfig{$61.0 \pm 1.0$} & \mainconfig{$18.1 \pm 0.3$} \\
    MeanFM     & 2 & $r$ & $0.933 \pm 0.041$ & $65.9 \pm 1.4$ & $27.3 \pm 0.2$ \\
                     &    & $c$ & $0.927 \pm 0.030$ & $98.0 \pm 10.4$ & $26.8 \pm 0.4$ \\
                     &    & $t$ & $0.993 \pm 0.009$ & $60.4 \pm 0.9$ & $27.1 \pm 0.3$ \\
    \midrule
    FM         & 1 & $r$ & $0.997 \pm 0.007$ & $68.3 \pm 0.9$ & $19.4 \pm 0.6$ \\
                     &    & \textbf{$c$} & \textbf{$0.997 \pm 0.007$} & \textbf{$67.5 \pm 1.2$} & \textbf{$18.9 \pm 0.2$} \\
                     &    & $t$ & $0.990 \pm 0.022$ & $67.1 \pm 1.6$ & $20.1 \pm 1.9$ \\
    FM         & 2 & $r$ & $0.997 \pm 0.007$ & $68.0 \pm 1.6$ & $27.6 \pm 0.5$ \\
                     &    & $c$ & $1.000 \pm 0.000$ & $68.4 \pm 3.0$ & $27.3 \pm 0.3$ \\
                     &    & $t$ & $0.983 \pm 0.020$ & $64.9 \pm 1.9$ & $27.9 \pm 0.5$ \\
    FM         & 20$^{\dagger}$ & $r$ & $0.960 \pm 0.070$ & $72.8 \pm 5.4$ & $496.5 \pm 74.8$ \\
                     &    & $c$ & $0.970 \pm 0.076$ & $65.6 \pm 4.3$ & $476.6 \pm 39.9$ \\
                     &    & $t$ & $0.987 \pm 0.032$ & $67.3 \pm 3.3$ & $465.8 \pm 36.0$ \\
    \midrule
    Diffusion  & 20$^{\dagger}$ & $r$ & $0.967 \pm 0.025$ & $73.9 \pm 8.7$ & $580.7 \pm 18.0$ \\
                     &    & \mainconfig{\baselinemark\ $c$} & \mainconfig{$0.983 \pm 0.025$} & \mainconfig{$69.0 \pm 11.5$} & \mainconfig{$563.5 \pm 29.6$} \\
                     &    & $t$ & $0.960 \pm 0.045$ & $82.8 \pm 14.7$ & $564.3 \pm 29.5$ \\
    \bottomrule
  \end{tabularx}
\end{table}
\dataref{analysis\_results\_checkpoint/19-09-UAV-Pillars-Exclude/batch\_avoiding\_combined\_20260919\_132703; DA\_in\_Paper/analysis/avoiding\_rules\_by\_protocol.py; agrees with DA\_20260827 \S10.1}
\begin{table}[htpb]
  \caption[What the twenty-episode evaluation covers]{What each cell of the twenty-episode evaluation
  rests on, counted from the corpus of record. A cell is complete at five training seeds on each of the
  three geometries; the entries below name every geometry that is short and the seeds it has. At
  $\nfe\in\{1,2\}$ every cell of every model is complete, and at $\nfe=20$ none is. These are also the
  slowest cells of the campaign, at half a second and more per control step against under thirty
  milliseconds.}
  \label{tab:coverage-k20}
  \centering
  \footnotesize
  \begin{tabularx}{\linewidth}{@{}l r L L@{}}
    \toprule
    Model & $\nfe$ & Projected arms ($r$, $c$, $t$) & Unprojected arm \\
    \midrule
    MeanFM, CI-MeanFM, FM, Diffusion & 1, 2 & complete & complete \\
    \midrule
    \mainconfig{\baselinemark\ Diffusion} & \mainconfig{20}
      & \mainconfig{\emph{both-hard}: seed 6 only under $r$ and $c$, absent under $t$}
      & \mainconfig{\emph{both-hard} absent} \\
    MeanFM & 20 & \emph{both-hard} absent; \emph{top-left-hard} 2 of 5 seeds & \emph{both-hard} absent; \emph{top-left-hard} 3 of 5 seeds \\
    FM & 20 & \emph{both-hard} 3 of 5 seeds under $r$ and $c$, 2 of 5 under $t$ & \emph{both-hard} 2 of 5 seeds \\
    \bottomrule
  \end{tabularx}
\end{table}
\dataref{seed-and-geometry census of
analysis\_results\_checkpoint/19-09-UAV-Pillars-Exclude/batch\_avoiding\_combined\_20260919\_132703;
DA\_in\_Paper/analysis/avoiding\_unprojected\_and\_budget20.py, block C}
\guard{The baseline row of \autoref{tab:state-headline} is therefore five seeds on two geometries plus
one seed on the third. The gap is recorded in
\texttt{data\_status/PENDING\_20260920\_all\_lacking\_runs.md} and is not closed here: the author has
ruled the re-run out on cost.}

\paragraph{Matched step budgets.}

\autoref{tab:avoiding-budget20} compares each flow-based model at two and twenty network evaluations
on the seed--geometry cells available at both budgets.

No model improves at twenty. Instantaneous-velocity matching goes from $1.000$ success with constraint
satisfaction to $0.970$, analytic average-velocity matching from $0.990$ to $0.955$, and
consistency-interpolated average-velocity matching from $1.000$ to $0.950$; the control-step column moves
by three steps in one direction for the first and by one to three in the other for the other two. Ten
times the network evaluations raise the time per control step for all three models.

The same ten-fold rise in the budget multiplies the time to
compute one action by $17$ for instantaneous-velocity matching and by about $36$ for both
average-velocity models, from under thirty milliseconds to between $0.48$ and $1.02$ seconds. At twenty
evaluations two of the three take longer per action than the diffusion baseline at the same budget.

\begin{table}[htpb]
  \caption[The flow-based models at their own budget and at the baseline's]{D3IL-avoiding at twenty
  episodes per seed, each flow-based model at $\nfe=2$ and at $\nfe=20$ under the selection rule it is
  operated with, against the baseline of record at its own budget. Tightened constraints, per-step
  projection, four candidate plans, temporal U-Net of $4.0$\,M parameters; mean $\pm$ the standard
  deviation over the seeds, aggregated as in \autoref{tab:state-headline}. \textbf{The two budgets of a
  model are restricted to the evaluation cells both of them cover}, so that the comparison is
  within-model and matched; the $\nfe=2$ numbers here can therefore differ slightly from
  \autoref{tab:state-headline}, which uses every cell a budget has. The last column states what each
  matched pair rests on.}
  \label{tab:avoiding-budget20}
  \centering
  \footnotesize
  \begin{tabularx}{\linewidth}{@{}l r c rrr L@{}}
    \toprule
    Model & $\nfe$ & Rule & S\&C & Steps & ms/step & Matched cells \\
    \midrule
    FM        & 2  & $c$ & $1.000 \pm 0.000$ & $68.8 \pm 2.8$  & $27.3 \pm 0.4$  & 5 seeds, 3 geometries; \emph{both-hard} 3 of 5 seeds \\
    FM        & 20 & $c$ & $0.970 \pm 0.076$ & $65.6 \pm 4.3$  & $476.6 \pm 39.9$ & \\
    \midrule
    MeanFM    & 2  & $t$ & $0.990 \pm 0.022$ & $61.7 \pm 2.6$  & $27.3 \pm 0.5$  & 5 seeds, 2 geometries; no \emph{both-hard} at $\nfe=20$, \emph{top-left-hard} 2 of 5 seeds \\
    MeanFM    & 20 & $t$ & $0.955 \pm 0.065$ & $63.5 \pm 10.0$ & $1003.9 \pm 43.9$ & \\
    \midrule
    CI-MeanFM & 2  & $t$ & $1.000$ & $59.7$ & $28.6$  & training seed 6, 3 geometries \\
    CI-MeanFM & 20 & $t$ & $0.950$ & $62.4$ & $1018.3$ & \\
    \midrule
    \mainconfig{\baselinemark\ Diffusion} & \mainconfig{20} & \mainconfig{$c$} & \mainconfig{$0.983 \pm 0.025$} & \mainconfig{$69.0 \pm 11.5$} & \mainconfig{$563.5 \pm 29.6$} & 5 seeds, 3 geometries; \emph{both-hard} seed 6 only \\
    \bottomrule
  \end{tabularx}
\end{table}
\dataref{analysis\_results\_checkpoint/19-09-UAV-Pillars-Exclude/batch\_avoiding\_combined\_20260919\_132703;
DA\_in\_Paper/analysis/avoiding\_unprojected\_and\_budget20.py, block B, matched cells}
\guard{Two limits. The consistency-interpolated rows are one training seed and carry no spread, so they
are read as a single run and not as a peer of the five-seed rows above them. And the $\nfe=20$
evaluations are the incomplete ones: analytic average-velocity matching has no \emph{both-hard}
evaluation at that budget at all, which is why its pair is matched on two geometries. The time column is
a measurement and is reported as one; this thesis does not explain why analytic average-velocity
matching at twenty evaluations costs more per action than the diffusion baseline at twenty, measured on
cells that are not the same.}

\FloatBarrier

\chapter{Reproducibility}\label{app:repro}

\section*{Quadrotor dimensions}

\begin{figure}[htpb]
  \centering
  \includegraphics[width=0.75\linewidth]{fig_platform_x2_dimensions}
  \caption[X2 dimensions and projection footprint]{The Skydio X2 rotor disks (radius $0.13$\,m), body
  and $0.54\times0.62$\,m collision footprint (dotted), drawn to scale. The dashed circle is the
  $0.31$\,m radius used to inflate the projection constraints.}
  \label{fig:platform-x2-dimensions}
\end{figure}
\FloatBarrier

\section*{Compute environment}

\autoref{tab:compute} is the machine of record for every number in this thesis. The hardware block
is transcribed from \texttt{sinfo}, \texttt{scontrol show node} and \texttt{lscpu} run on the cluster
on 2026-09-09, cross-checked against the provenance header (\texttt{NODE}, \texttt{GPU INFO},
\texttt{GIT REV}) that every batch job in this project writes to its own log. The software block is
transcribed from the environment-verification job that certified the cluster installation, and from
the pinned dependency list.

\begin{table}[htpb]
  \caption[Compute environment]{The machine and software stack of record. All runs used a single
  node; every job was allocated one GPU.}
  \label{tab:compute}
  \centering
  \begin{tabularx}{\linewidth}{@{}L L@{}}
    \toprule
    Component & Configuration \\
    \midrule
    \multicolumn{2}{l}{\emph{hardware}}\\
    Node                    & single shared GPU node, institute cluster \\
    CPU                     & $2\times$ AMD EPYC 7282, 16 cores each \\
    Cores / threads         & 32 physical / 64 logical, 2 NUMA domains, 1.5--2.8\,GHz \\
    Cache                   & 512\,KiB L2 per core; 128\,MiB L3 total ($8\times16$\,MiB) \\
    Main memory             & 252\,GiB; no node-local scratch \\
    GPUs on node            & $8\times$ NVIDIA RTX A5000, 24\,GB each \\
    GPUs per job            & \textbf{1}, held exclusively (no oversubscription, no preemption) \\
    NVIDIA driver           & 530.41.03, unchanged over the project \\
    \midrule
    \multicolumn{2}{l}{\emph{scheduling}}\\
    Scheduler               & SLURM 21.08.5 \\
    Wall-clock limit        & 24\,h per job, partition-enforced \\
    Typical training job    & 1 GPU, 8 threads, 32\,GB, 3--24\,h \\
    Operating system        & Ubuntu, Linux 5.15.0 \\
    \midrule
    \multicolumn{2}{l}{\emph{software}}\\
    Python                  & 3.10 (Conda environment) \\
    PyTorch                 & 2.2.2, CUDA 12.1 build \\
    NumPy / SciPy           & 1.26.4 / 1.13.1 \\
    MuJoCo                  & 2.3.7, headless EGL rendering \\
    Projection solver       & SciPy SLSQP (default); CasADi 3.6.5 with IPOPT selectable \\
    ODE integration         & \texttt{torchdiffeq} 0.2.3 \\
    Diffusers / Transformers & 0.31.0 / 4.41.2 \\
    Gym / Gymnasium         & 0.26.2 / 0.29.1 \\
    \bottomrule
  \end{tabularx}
\end{table}

Two qualifications belong with the table. The software block is a snapshot taken when the
environment was certified at the start of the project; no version change is recorded in any
subsequent job log, but the environment was not frozen, so the defensible claim is that these
versions were verified at the outset rather than that each one produced every reported number.
And the accelerator, though the GPU is held exclusively, sits in a node whose CPU and memory are
shared, which is why \autoref{sec:setup:protocol} declines to read the timing figures as a hardware
benchmark.

Two properties of this machine bear on how the times per action of \autoref{ch:results} are read. The
node is shared: it serves several GPU queues at once, but the scheduler neither oversubscribes GPUs nor
preempts running jobs, so a job holds its accelerator exclusively and contention is confined to the CPU,
the memory system and I/O. And thirty-two physical cores serve eight accelerators, while both the MuJoCo
rollout and the per-step projection, solved by SLSQP (\autoref{sec:method:proj}), run on the CPU, so
evaluation throughput is plausibly limited by the CPU rather than the RTX~A5000. Times are therefore
compared between configurations measured on this same node, and not read as hardware benchmarks.

Reproducing a run requires the batch scripts as much as the code: the scheduling parameters, the
GPU pinning and the headless-rendering configuration all live in the job scripts, and each job
records the git revision it ran at in its own log header.
\srcnote{Slurm\_Codes/submit.sh; Slurm\_Codes/sbatch/}

\hole{Total compute spent is not yet quantified. Arithmetic over the batch-log timestamps gives an
order of $10^{3}$ GPU-hours across roughly five months, but that covers only the jobs whose logs
record both a start and an end; re-derive it with \texttt{sacct} before stating any figure.}


% -----------------------------------------------------------------------------
%  v3 ADDITION. Two tables the results chapters create a need for:
%  (1) the reverse name map -- Chapters 5-7 use mechanism names exclusively
%      (Auxiliary/Naming/NAMING_20260910_master_table.md rule 2), so without
%      this table a reader cannot match a results row to a released checkpoint
%      or log. NAMING section 8 records that debt; this discharges its first half.
%  (2) the corpora of record, so a number can be traced to a batch.
% -----------------------------------------------------------------------------

\section*{Names as they appear in the released artefacts}

This thesis names methods by their mechanism and never by the tokens that identify them in
configuration files, checkpoint paths and evaluation logs. \autoref{tab:names} maps the names used in \autoref{ch:results} to those tokens.

\begin{table}[htpb]
  \caption[Thesis names and released artefact tokens]{Thesis name against the token a reader will
  meet in the released configuration files, checkpoint directories and evaluation logs.}
  \label{tab:names}
  \centering
  \begin{tabularx}{\linewidth}{@{}L L@{}}
    \toprule
    Name used in this thesis & Token in the artefacts \\
    \midrule
    \multicolumn{2}{l}{\emph{generative models}}\\
    Diffusion & \texttt{GaussianDiffusion}, \texttt{ddpm}, \texttt{diffusion} \\
    FM        & \texttt{FlowMatchingODE}, \texttt{fm}, \texttt{flow\_matcher\_v3} \\
    MeanFM    & \texttt{MeanFlowODE}, \texttt{mf} \\
    CI-MeanFM & \texttt{AlphaFlowODE}, \texttt{af}; $\alpha_{\mathrm{end}}$ as \texttt{ae0.2} / \texttt{ae0.05} \\
    \midrule
    \multicolumn{2}{l}{\emph{projection methods}}\\
    unguided (no projection) & the \texttt{diffuser} variant \\
    per-step projection (DPCC) & \texttt{dpcc-r} / \texttt{-c} / \texttt{-t} \\
    endpoint projection & \texttt{hardflow\_sls-*}, \texttt{hardflow\_new-*} \\
    tightened constraints & the \texttt{-tightened} suffix \\
    \midrule
    \multicolumn{2}{l}{\emph{selection rules}}\\
    random                              & the \texttt{-r} suffix \\
    cumulative projection cost          & the \texttt{-c} suffix \\
    temporal consistency                & the \texttt{-t} suffix \\
    \midrule
    \multicolumn{2}{l}{\emph{parameters and ablations}}\\
    step budget $\nfe$                  & \texttt{K}, \texttt{n\_diffusion\_steps}, \texttt{flow\_steps} \\
    activation threshold $\eta$         & \texttt{A}, \texttt{activation\_threshold} \\
    guiding steps                       & \texttt{n\_genuine} \\
    corrected scene geometry            & \texttt{u7hg}, ``honest geometry'' \\
    constraint-family ablations         & \texttt{geo\_free}, \texttt{bounds\_free}, \texttt{model\_free} \\
    \bottomrule
  \end{tabularx}
\end{table}

Two tokens name something other than what the thesis means by the word. \texttt{diffuser} denotes the ancestor codebase, the published method of the same name, and the variant evaluated without projection. The conditioning flag \texttt{filmv1} selects the feature-wise conditional biasing of \autoref{sec:method:backbone}, which is what every reported visual result uses.

\section*{Corpora of record}

\autoref{tab:corpora} lists the evaluation batches every number in \autoref{ch:results} was computed
from, with the protocol that produced each. The figure scripts read these directories directly. The two
batches of 2026-09-19 are the corpora of record: they supersede their predecessors and reproduce every
number those had already produced, so what they add are rows and not revisions. They carry the evaluations
of UAV-pillars as well, which are excluded from the analysis and reported nowhere
(\autoref{sec:res:uav}).

\begin{table}[htpb]
  \caption[Corpora of record]{The evaluation batches behind the results, and their protocols.}
  \label{tab:corpora}
  \centering
  \small
  \begin{tabularx}{\linewidth}{@{}L L L@{}}
    \toprule
    Environment and results & Batch & Protocol \\
    \midrule
    D3IL-avoiding, all protocol tables & \texttt{batch\_avoiding\_combined\_20260919} & 5 seeds; 2 and 20 episodes per geometry \\
    D3IL-avoiding, generative models & \texttt{batch\_avoiding\_combined\_20260825} & 5 seeds $\times$ 20 episodes \\
    D3IL-avoiding, CI-MeanFM & \texttt{batch\_avoiding\_combined\_20260903} & seed 6, 20 episodes \\
    D3IL-avoiding, projection methods & \texttt{batch\_avoiding\_combined\_20260906} & 4 seeds $\times$ 3 geometries $\times$ 2 episodes \\
    D3IL-aligning & \texttt{batch\_va2\_20260915} & seed 6, 10 paired contexts \\
    UAV-corridor and UAV-s-curve & \texttt{batch\_uav\_20260919} & seed 6, 12 flights (4 per route) and 10 flights \\
    UAV-corridor, earlier reference & \texttt{batch\_uav\_20260914} & seed 6, 12 flights (4 per route) \\
    \bottomrule
  \end{tabularx}
\end{table}

\hole{Complete \autoref{tab:corpora} with the SLURM job identifiers and the git revision each batch
ran at; both are recorded in every job log header and in the project's run ledger. Add the
checkpoint-tag spellings, which are the other half of what \autoref{tab:names} owes a reader.}
